% Generated draft: requires mathematical and source review.
\documentclass[11pt]{article}
\usepackage[margin=25mm]{geometry}
\usepackage{fontspec}
\setmainfont{Latin Modern Roman}
\usepackage{amsmath,amssymb,mathtools,unicode-math}
\setmathfont{Latin Modern Math}
\usepackage{enumitem,needspace,xurl,hyperref}
\setlength{\parindent}{0pt}
\setlength{\parskip}{5pt}
\setlength{\emergencystretch}{3em}
\newcommand{\N}{\mathbb{N}}
\newcommand{\Z}{\mathbb{Z}}
\newcommand{\Q}{\mathbb{Q}}
\newcommand{\R}{\mathbb{R}}
\newcommand{\C}{\mathbb{C}}
\newcommand{\F}{\mathbb{F}}
\newcommand{\A}{\mathbb{A}}
\newcommand{\PP}{\mathbb{P}}
\newcommand{\E}{\mathbb{E}}
\newcommand{\eps}{\varepsilon}
\newcommand{\dd}{\,\mathrm{d}}
\begin{document}
\section*{UnsolvedMath ID: 801190}
\textbf{Source problem code:} catalogue entry
\subsection{Definitions and mathematical statement}
Let \(n\geq 1\) be an integer, and let \(\lambda\vdash n\) mean that \(\lambda=(\lambda_1,\ldots,\lambda_\ell)\) is a weakly decreasing sequence of positive integers with sum \(n\). Let \(S_n\) be the symmetric group on \(\{1,\ldots,n\}\). For \(w\in S_n\), write \(\rho(w)\) for the partition consisting of its cycle lengths. In commuting indeterminates \(x_1,x_2,\ldots\), put
\[
p_r=\sum_{i\geq 1}x_i^r,\qquad p_\mu=\prod_{j=1}^{\ell(\mu)}p_{\mu_j}.
\]
Let \(m_\lambda\) be the monomial symmetric function: the sum of all distinct monomials whose positive exponents, arranged in decreasing order, are the parts of \(\lambda\).

A class function on \(S_n\) is a function constant on conjugacy classes. Its Frobenius characteristic is
\[
\operatorname{ch}(f)=\frac{1}{n!}\sum_{w\in S_n}f(w)p_{\rho(w)}.
\]
Define \(\phi^\lambda\) to be the unique class function satisfying \(\operatorname{ch}(\phi^\lambda)=m_\lambda\). For a real matrix \(A=(a_{ij})_{1\leq i,j\leq n}\), define
\[
\operatorname{Imm}_{\phi^\lambda}(A)
=\sum_{w\in S_n}\phi^\lambda(w)\prod_{i=1}^{n}a_{i,w(i)}.
\]
The matrix \(A\) is \emph{totally nonnegative} if, for every \(1\leq k\leq n\) and all increasing index sequences \(i_1<\cdots<i_k\) and \(j_1<\cdots<j_k\),
\[
\det\bigl(a_{i_r,j_s}\bigr)_{r,s=1}^{k}\geq 0.
\]
Zero minors and singular matrices are allowed.

\textbf{Conjecture (attributed in the supplied record to John R. Stembridge; Stanley's Problem 17).} For every integer \(n\geq1\), every partition \(\lambda\vdash n\), and every totally nonnegative matrix \(A\in\R^{n\times n}\),
\[
\operatorname{Imm}_{\phi^\lambda}(A)\geq0.
\]
The supplied registry labels this conjecture open; that label is supplied curation and is not an independently verified assessment of its current status.
\subsection{Short English statement}
Stembridge's conjecture asks whether every immanant associated with a monomial symmetric function is nonnegative on every real totally nonnegative square matrix. The claim includes matrices with zero minors. It is listed as Problem 17 in Stanley's 1999 collection and marked open in the supplied registry.
\subsection{Sources}
\begin{itemize}
\item Supplied local UnsolvedMath registry record: permanent ID \(801190\), exact source problem code \texttt{catalogue entry}, in the AMR Open Problem Lists collection. The record lists Stanley as proposer and 1999 as the proposed year, while its background attributes the conjecture to John R. Stembridge. Its supplied status is \texttt{open}.
\item Stanley, \emph{Positivity Problems and Conjectures in Algebraic Combinatorics} (1999), Problem 17. Supplied source URL: \url{https://math.mit.edu/~rstan/papers/problems.pdf}. The record gives an access date of 2026-10-16 and an extraction designation of \texttt{pdf+html\_verified}. These are supplied dated provenance and verification annotations; no independent source visit or verification is claimed here.
\item John R. Stembridge, \emph{Bulletin of the London Mathematical Society} \textbf{23} (1991), 422--428. This citation and the attribution of the conjecture are supplied in the record's background; no article title or direct article URL is supplied.
\end{itemize}
\end{document}
